\section{Label Validation}
\label{app:judge}
\paragraph{Judges.}
All judges receive the same system prompt and item prompt (Appendix~\ref{app:prompts}): the question, the presupposition stated explicitly, and the response truncated to 1{,}200 characters, with four options (accept, rebut, cannot verify, refuse). Claude Haiku 5.5 (primary) is called through the Anthropic API with thinking disabled, low effort and at most 16 output tokens; a label is accepted only from a reply that is the letter itself or an explicit ``Label: X'', and otherwise the item is asked again (\texttt{code/judge\_api.py}). Qwen2.5-7B-Instruct and Mistral-7B-Instruct-v0.3 run locally and take the option with the highest next-token log-probability. Every judge labels every response. Refusals are merged into \emph{cannot verify}; the 3-judge majority breaks three-way ties as \emph{cannot verify}.

\paragraph{First validation sample.}
Claude Sonnet 5.5 labelled a stratified random sample of main-grid responses (seed 0; up to 8 responses per model $\times$ condition $\times$ item type, knowledgeable models; \texttt{code/silver\_labels.py}) with the same prompt, allowed to reason briefly before giving its label. We chose the primary judge on this sample (Table~\ref{tab:silver}). The table reports label accuracy (acc.), Cohen's $\kappa$, and rebuttal precision (P) and recall (R) against Sonnet's automatic reference labels, called silver labels. Haiku combines high rebuttal recall with the highest overall agreement; these references are not human judgements. Mistral-7B misses many rebuttals, mostly labelling graded corrections (``as of my last update, $Z$ was the CEO'') as \emph{cannot verify}.

\paragraph{Held-out validation.}
To assess the choice on a separate sample, Claude Sonnet 5.5 (letter only) and OpenAI gpt-5.3-codex (low reasoning effort; \texttt{code/judge\_openai.py}) each labelled every response to the entity triples under \textsc{Plain}, \textsc{Premise-check} and \textsc{Cutoff-stated} (2{,}205 responses) and, in round three, \textsc{RAG-fresh+PC} (735), for all seven models. For \WSilRetry{} of the \WHoKeys{} items Sonnet's letter-only reply began a sentence that the 8-token cap cut off; they were re-asked with a short reasoning budget ending in ``Label: X'', and \WSilMissing{} items remain unlabelled. The re-asked items are harder than average, so including them lowers agreement slightly. Sonnet and Haiku are both Claude models; gpt-5.3-codex is the only labeller from another family. Table~\ref{tab:heldout} reports accuracy and Cohen's $\kappa$ for the three initial conditions, followed by \textsc{RAG-fresh+PC}. The consensus rows use only responses on which Sonnet and gpt-5.3-codex agree; the ``all'' rows include every available label. None of these held-out labels was used to choose the primary judge.

Table~\ref{tab:judges3} recomputes the main rates and margins under each labeller. Haiku, Qwen and Mistral denote the three judges above; Maj.\ is their majority vote. Sonnet and codex denote the held-out labellers. PC, Cutoff and RAG-f+PC abbreviate \textsc{Premise-check}, \textsc{Cutoff-stated} and \textsc{RAG-fresh+PC}. The former-holder margin depends most on the labeller: Mistral finds none under \textsc{Plain}, whereas both held-out labellers retain a positive margin.

\begin{table}[t]\centering\small\setlength\tabcolsep{4pt}
\begin{tabular}{@{}l ccc@{}}
\toprule
Comparison & $n$ & acc. & $\kappa$ \\
\midrule
Sonnet 5.5 vs.\ gpt-5.3-codex & 2204 & 0.95 & 0.92 \\
\midrule
Haiku 5.5 vs.\ consensus & 2090 & 0.94 & 0.91 \\
3-judge majority vs.\ consensus & 2090 & 0.95 & 0.92 \\
Qwen2.5-7B vs.\ consensus & 2090 & 0.91 & 0.86 \\
Mistral-7B vs.\ consensus & 2090 & 0.84 & 0.76 \\
\midrule
Haiku vs.\ Sonnet (all) & 2204 & 0.91 & 0.86 \\
Haiku vs.\ gpt-5.3-codex (all) & 2205 & 0.92 & 0.88 \\
RAG-f+PC: Sonnet vs.\ codex & \hphantom{0}734 & 0.93 & 0.85 \\
RAG-f+PC: Haiku vs.\ Sonnet & \hphantom{0}734 & 0.90 & 0.77 \\
RAG-f+PC: Haiku vs.\ codex & \hphantom{0}735 & 0.87 & 0.71 \\
\bottomrule\end{tabular}
\caption{Held-out label agreement on entity triples. Consensus uses only responses where both held-out labellers agree; \emph{All} uses every available label. The last three rows assess \textsc{RAG-fresh+PC}.}
\label{tab:heldout}\end{table}

\begin{table}[t]\centering\small\setlength\tabcolsep{4pt}
\begin{tabular}{@{}l rrrrc@{}}
\toprule
 & & & \multicolumn{3}{c}{rebut} \\ \cmidrule(l){4-6}
Judge & acc. & $\kappa$ & P & R & rate (\%) \\
\midrule
Haiku 5.5 (primary) & 0.82 & 0.73 & 0.91 & 0.95 & 42 \\
3-judge majority & 0.81 & 0.72 & 0.94 & 0.91 & 39 \\
Qwen2.5-7B & 0.80 & 0.70 & 0.88 & 0.92 & 42 \\
Mistral-7B & 0.70 & 0.56 & 0.99 & 0.64 & 26 \\
\midrule
Silver (Sonnet 5.5) & \multicolumn{4}{c}{--} & 40 \\
\bottomrule\end{tabular}
\caption{Agreement with automatic Sonnet 5.5 labels (602 responses). P/R: rebuttal precision/recall. These reference labels are not human judgements.}
\label{tab:silver}\end{table}

\begin{table}[t]\centering\small\setlength\tabcolsep{3.1pt}
\begin{tabular}{@{}l r rrr rr@{}}
\toprule
 & \multicolumn{1}{c}{primary} & \multicolumn{3}{c}{other judges} & \multicolumn{2}{c}{held-out} \\ \cmidrule(lr){2-2}\cmidrule(lr){3-5}\cmidrule(l){6-7}
 & Haiku & Maj. & Qwen & Mistral & Sonnet & codex \\
\midrule
\multicolumn{7}{@{}l}{\textit{rebut under Plain (\%)}} \\
\quad \NEW{} & 30 & 30 & 33 & 23 & 31 & 30 \\
\quad \ANC{} & 41 & 40 & 45 & 20 & 42 & 41 \\
\quad \SWP{} & 40 & 39 & 41 & 33 & 38 & 39 \\
\multicolumn{7}{@{}l}{\textit{$J_{\text{fam}}$ (pts)}} \\
\quad Plain & +11 & +11 & +12 & $-$2 & +11 & +11 \\
\quad PC & +13 & +10 & +11 & +8 & +9 & +10 \\
\quad Cutoff & +12 & +12 & +13 & +4 & +16 & +16 \\
\multicolumn{7}{@{}l}{\textit{$J_{\text{unf}}$ (pts)}} \\
\quad Plain & +10 & +9 & +9 & +10 & +7 & +9 \\
\quad PC & +13 & +12 & +12 & +10 & +10 & +12 \\
\quad Cutoff & +8 & +7 & +9 & +7 & +7 & +7 \\
\multicolumn{7}{@{}l}{\textit{rebut under RAG-f+PC (\%)}} \\
\quad \NEW{} & 39 & 40 & 54 & 37 & 27 & 28 \\
\bottomrule\end{tabular}
\caption{Entity-triple results under each labeller (seven models, same responses). Maj.: majority vote; PC: our premise check. The text defines the remaining abbreviations.}
\label{tab:judges3}\end{table}

